\section{環境管理原則と社会の普遍的目標}
\label{theory:env:principles}
\label{theory:env:sdgs}
汚染者負担原則は汚染防止・対策費用を誰が負担するか、拡大生産者責任は製品の使用後まで生産者の責任をどう広げるかという異なる視点である。回収・再利用までの責任は金銭的・物理的な制度設計を含み、一律に製造者だけが全実務を担当するという意味ではない。\cite{env59_epr}

2030アジェンダのSDGsは全ての国を対象とし、環境・社会・経済を統合する17ゴール・169ターゲットから構成される。途上国限定とすると普遍性を落とす。国内実施体制や数値を定める際も、目標とその法的強制手段、ゴールとターゲットを区別する。\cite{env_sdgs_un,env_sdgs2016}

\section{緩和の目標、適応と生態系の機能}
\label{theory:env:paris}
\label{theory:env:ecodrr}
NDCの提出は温室効果ガスの排出削減等の緩和に対応し、五年ごとに後続目標を提出する仕組みである。協定の温度目標、透明性、国際協力、森林の取組は別条項として読む。削減の目標があることを「影響に備える適応」という説明へ置き換えない。\cite{env59_ndc}

Eco-DRRは湿原の遊水、森林の保全等の生態系機能を防災・減災に生かす。適応と生物多様性に共便益があり、地域特性と複数分野の連携も必要になる。自然を利用するという名称だけで全災害の安全を保証しない。\cite{env59_ecodrr}

\section{一次供給、最終消費、人口・GDPの原単位}
\label{theory:env:home}
\label{theory:env:oil2017}
\label{theory:env:primaryshare}
\label{theory:env:energytrends}
一次エネルギー供給は発電用の投入等を含み、最終消費は利用部門への供給を表す。一次供給に占める再エネ比と、発電量に占める再エネ比は違う。未活用エネルギーを含めるという問題独自の境界があれば、その指定を使う。

家庭の総エネルギーは電気だけでなく燃料を含む。用途、世帯数、一人当たり、機器効率、気候の要因を分ける。総量が減っても全世帯で毎年減るとは限らない。GDP当たりの比較には同じエネルギー範囲と通貨・実質化方式を用いる。原油の数量依存度と総輸入額比も分母が違う。出題後の統計遡及改訂は当時公表値と別に記録する。\cite{env59_energy2019}

\section{炭素換算から二酸化炭素換算へ進む}
\label{theory:env:nationalco2}
燃料投入を炭素排出係数で重み付けする。係数がt-C/TJならPJを1000倍してTJへ直し、燃焼に用いない原料分を先に除く。
\[
 M_{\mathrm{CO_2}}=\frac{44}{12}\sum_i E_i e_i .
\]
係数が既にCO\(_2\)換算であればさらに \(44/12\) を掛けない。投入熱量全体、化石燃料分、非エネルギー分を同じ総量として掛けると境界が変わる。最終値をMtへ直して候補の桁と比較する。

\section{家庭の計測と管理を別の機能として読む}
\label{theory:env:smart}
\label{theory:env:airhistory}
通信機能を持つスマートメーターは遠隔検針・データ取得等の計測装置である。HEMSは機器の状態やデータを使う家庭のエネルギー管理であり、地域の管理、系統全体の需給管理とは範囲が違う。計測装置を管理・最適化システムそのものと定義しない。\cite{env59_smart}

汚染の歴史を読む際も対象を分ける。工業型の燃料硫黄、自動車の窒素酸化物と粒子、日射による光化学オキシダントの二次生成はそれぞれ発生源と反応が違う。規制や監視の導入後に改善しても全物質・全地域が同じ達成率になるとはいえない。

\section{発見と応用は一方向の年表ではない}
\label{theory:env:historydiscover}
\label{theory:env:technologysequence}
\label{theory:env:japanindustrial}
経験的な技術が機序の科学的解明より先に成立することがある。Jennerの種痘は1796年で、ウイルスの発見より前だった。一方、装置が観測・発見を可能にし、理論が新たな素子開発を導く例もある。核分裂と発電、放射線と加速器、電子管と半導体は各段階を分ける。\cite{env59_jenner,env59_lawrence}

日本の技術史では、科学的管理法、統計的手法の個別利用、統計的品質管理の組織的普及を同じ導入年にしない。1950年の日科技連によるDeming講義は普及の節目であって、戦前に一例も統計利用がなかったという主張ではない。学協会の高学歴者構成と欧米との差も組織・時期による。令和2年度の公式追加正答3と4は、その比較表現の留保を示す。\cite{env59_juse}

\section{リスクの評価・管理・知覚・対話}
\label{theory:env:risk}
科学的リスク評価は影響や可能性を不確実性込みで推定する。管理は技術可能性、費用、公平性などを含む判断である。知覚は経験、信頼、自発性等にも依存するため評価と一致するとは限らない。対話は一方向の説得ではなく、情報と意見の相互交換となる。科学的根拠による規制判断の支援と、社会の価値判断を科学だけで確定する主張を分ける。\cite{eth54_riskcomm}

\section{プラスチック循環は粒径、流出、原料を分ける}
\label{theory:env:plastics}
マイクロプラスチックは概ね5mm以下という大きさ、海洋への流出は廃物管理と陸域・河川の経路、再生利用は処理と原料の区分から考える。熱回収を材料再生と同じ量へ二重計上しない。3R+Renewableは利用抑制・再使用・再生利用に再生可能資源への代替を加える。

政府の資源循環戦略策定日は2019年5月31日。令和2年度原典の6月という表現とは一致しないため、内容の正誤と日付の厳密な判定を分け、公式正答と疑義をともに残す。バイオ由来と生分解性は異なる性質で、全て海中ですぐ無害化するという意味ではない。\cite{env59_plastic}
